% =====================================================================
%  Seance 5 -- B. detecter par ML : sessions, features, IsolationForest
% =====================================================================
\section{Détecter par apprentissage automatique}

\begin{frame}{Étape 1 : transformer 100 000 lignes en sessions}
\begin{columns}[t]
\begin{column}{0.52\textwidth}
\begin{kbox}[Session = (machine, fenêtre de 10 min)]
\begin{itemize}
\item \py{fenetre = horodatage.dt.floor("10min")} puis
\py{groupby(["machine","fenetre"])}.
\item 100\,000 lignes $\rightarrow$ \textbf{10\,000 sessions}
(10 machines $\times$ 1\,008 fenêtres sur 7 j).
\item Une session est la bonne granularité : ni l'événement (trop bruité),
ni le serveur-jour (trop grossier).
\end{itemize}
\end{kbox}
\begin{ebox}[Chiffres mesurés sur nos logs]
\begin{itemize}
\item 4\,828 lignes en échec (\py{auth\_reussie = False}) sur 7 jours,
soit $\approx$ \textbf{690 / jour} sur le parc --- la mesure corrige
l'ordre de grandeur « $\sim$30 / jour » évoqué en séance 1, qui ne comptait
que certains messages d'erreur.
\item \textbf{12 sessions d'attaque} (bdd-principal, 02 h -- 04 h le 05/10)
sur 10\,000 : \textbf{0,12 \%} des sessions. Chercher une aiguille
dans 10\,000 bottes de foin : c'est le problème de l'analyste SOC.
\end{itemize}
\end{ebox}
\end{column}
\begin{column}{0.44\textwidth}
\begin{center}
\includegraphics[width=\textwidth]{f29_echecs_7j.png}
\end{center}
\end{column}
\end{columns}
\end{frame}

\begin{frame}{Étape 2 : les features d'attaque par session}
\begin{center}
\footnotesize
\begin{tabular}{@{}p{4.1cm}p{5.4cm}p{2.6cm}@{}}
\toprule
\textbf{Feature} & \textbf{Intuition de détection} & \textbf{Attaque mesurée}\\
\midrule
\py{echecs} & un brute force empile les échecs & moyenne \textbf{10,4}, max 15\\
\py{taux\_echec} & échecs / événements de la session & $= 1{,}0$ (zéro succès)\\
\py{echecs\_consecutifs} & rafale ininterrompue $=$ signature forte & jusqu'à 15\\
\py{utilisateurs} & le robot teste des comptes en série & 3,7 en moyenne\\
\py{entropie} & diversité des comptes testés (bits) & 1,73 vs 2,22 en normal\\
\py{comptes\_sensibles} & échecs sur \py{root}, \py{admin}, \py{inconnu-*} & \textbf{100 \%} des échecs\\
\py{nuit} & fenêtre 22 h -- 06 h & toutes les sessions\\
\bottomrule
\end{tabular}
\end{center}
\vspace{1mm}
\begin{kbox}[Aucune feature ne suffit seule]
Le maximum de la baseline est \textbf{4 échecs / session} --- un seuil sur
\py{echecs} marcherait ici. Mais une session normale de nuit à 3 échecs sur
des comptes variés a une \textbf{entropie haute} (2,2 bits) alors que
l'attaque teste une \emph{petite liste} de comptes (1,73 bit). Combiner
$\ge 2$ features faibles $=$ l'idée centrale du ML de sécurité.
\end{kbox}
\end{frame}

\begin{frame}{L'entropie de Shannon des comptes testés}
\begin{columns}[t]
\begin{column}{0.5\textwidth}
\begin{block}{La formule (en bits)}
$H = -\sum_i p_i \log_2 p_i$, où $p_i$ est la part des essais portant sur
le compte $i$.
\begin{itemize}
\item Tous les essais sur \emph{un seul} compte : $H = 0$.
\item Essais uniformément répartis sur $k$ comptes : $H = \log_2 k$.
\item Un brute force humain outillé teste une \textbf{petite liste ordonnée}
(\py{root}, \py{admin}, comptes par défaut) : entropie \emph{basse} à
l'échelle de la session.
\end{itemize}
\end{block}
\begin{ebox}[En pratique avec scipy]
\py{scipy.stats.entropy(comptage, base=2)} --- ou à la main :
\py{-sum(p*log2(p) for p in frequences)}. C'est la tâche B2 du TP.
\end{ebox}
\end{column}
\begin{column}{0.46\textwidth}
\begin{center}
\includegraphics[width=\textwidth]{f30_entropie_echecs.png}
\end{center}
\end{column}
\end{columns}
\end{frame}

\begin{frame}[fragile]{IsolationForest : isoler l'anormal, pas décrire le normal}
\begin{columns}[t]
\begin{column}{0.52\textwidth}
\begin{block}{Le principe en trois phrases}
Une forêt d'arbres coupe l'espace des features au hasard. Une session
\textbf{isolée} (loin des autres, dans une région clairsemée) tombe dans une
feuille après \textbf{peu de coupes} ; une session du flot normal demande
beaucoup de coupes. Le score d'anomalie mesure cette facilité d'isolement.
\end{block}
\begin{itemize}
\item Non supervisé : \textbf{aucune étiquette} d'attaque nécessaire.
\item \py{contamination} : la fraction d'anomalies attendue (ici 0,005)
--- c'est \emph{notre a priori}, pas une vérité mesurée.
\item Alternatives : \py{LocalOutlierFactor} (densité locale),
autoencodeurs (reconstruction).
\end{itemize}
\end{column}
\begin{column}{0.44\textwidth}
\begin{codebox}[Le cœur de B3]
\begin{lstlisting}
iso = IsolationForest(
    n_estimators=200,
    contamination=0.005,
    random_state=5534, n_jobs=2)
iso.fit(X)                    # X = sessions x 7 features
score = -iso.decision_function(X)
top = S.assign(score=score)\
       .sort_values("score",
                    ascending=False)
top30 = top.head(30)          # l'analyste
                              # ne lit que ca
\end{lstlisting}
\end{codebox}
\end{column}
\end{columns}
\vspace{1mm}
\begin{center}\small
Seuil de décision $\ne$ vérité : on juge le détecteur par son
\textbf{classement} (précision@N), pas par une étiquette binaire.
\end{center}
\end{frame}

\begin{frame}{Résultat réel : toutes les sessions d'attaque au top 12}
\begin{columns}[t]
\begin{column}{0.5\textwidth}
\begin{center}
\includegraphics[width=\textwidth]{f31_precision_at_n.png}
\end{center}
\end{column}
\begin{column}{0.46\textwidth}
\begin{ebox}[Chiffres mesurés (graine 5534)]
\begin{itemize}
\item Rangs 1 à 12 : \textbf{exactement les 12 sessions d'attaque}.
\item \textbf{précision@10 $= 1{,}00$}, précision@30 $= 0{,}40$,
précision@50 $= 0{,}24$.
\item Rappel (top 30) $= 1{,}00$ : rien n'est manqué.
\item Par hasard, précision@30 serait $\approx$ 0,0012 :
\textbf{$\times$ 333 mieux que le hasard}.
\end{itemize}
\end{ebox}
\begin{wbox}[Lire la décroissance]
Au-delà du rang 12, précision@N décroît : les sessions « anormales mais
normales » (4 échecs de service, 7 utilisateurs) remontent. \textbf{Le
classement est un curseur} : l'analyste descend jusqu'à ce que le coût de
lecture dépasse la valeur des alertes.
\end{wbox}
\end{column}
\end{columns}
\end{frame}

\begin{frame}{Regrouper les alertes : le clustering rattache l'attaque}
\begin{columns}[t]
\begin{column}{0.5\textwidth}
\begin{block}{Pourquoi regrouper ?}
12 sessions d'attaque $=$ 12 alertes séparées... ou \textbf{1 incident} sur
2 heures. Sans regroupement, l'analyste croit à 12 problèmes distincts ;
avec, il voit \emph{une} campagne coordonnée sur \py{bdd-principal}.
\end{block}
\begin{itemize}
\item \py{DBSCAN} sur (machine encodée, heure, features) : les sessions
d'attaque partagent machine et fenêtres adjacentes $\Rightarrow$ un seul
cluster dense, le bruit reste \py{-1}.
\item Plus simple et robuste : \textbf{fusionner} les sessions anormales de
la même machine séparées de $\le$ 2 fenêtres.
\end{itemize}
\end{column}
\begin{column}{0.46\textwidth}
\begin{kbox}[De session à incident]
La sortie du détecteur n'est pas « 12 alertes » mais :
\begin{center}
\emph{« Campagne de brute force,\\ bdd-principal, 05/10 02 h -- 04 h,\\
125 échecs, comptes root/admin/inconnu-* »}
\end{center}
C'est la granularité que l'exploitant sait traiter --- et la base du
\textbf{livrable B5} du TP.
\end{kbox}
\end{column}
\end{columns}
\end{frame}

\begin{frame}{Quiz 2 --- sessions et détection non supervisée (5 min)}
\begin{qbox}[3 questions]
\begin{enumerate}
\item Pourquoi agréger en sessions de 10 min plutôt qu'analyser ligne par ligne ?\\
(a) c'est plus rapide \quad (b) une ligne ne dit rien : l'anomalie est dans la
\emph{densité} et la \emph{combinaison} d'événements \quad (c) parce que les
logs sont trop gros
\item Une session teste 8 comptes parfaitement équilibrés. Entropie :\\
(a) 0 bit \quad (b) 3 bits \quad (c) 8 bits
\item \py{contamination=0.005} dans IsolationForest signifie :\\
(a) 0,5 \% des données sont infectées \quad (b) un a priori sur la fraction
d'anomalies, qui calibre le seuil de décision \quad (c) le taux de faux positifs
\end{enumerate}
\end{qbox}
\vspace{1mm}
\begin{solbox}[Réponses]
\textbf{1 : b} --- l'agrégation fabrique les features qui portent la signature.
\quad
\textbf{2 : b} --- $\log_2 8 = 3$ bits, distribution uniforme.
\quad
\textbf{3 : b} --- c'est un paramètre d'a priori, à ajuster selon le contexte.
\end{solbox}
\end{frame}
